\documentclass[11pt,a4paper]{report}

% --- Packages ---
\usepackage[utf8]{inputenc}
\usepackage[T1]{fontenc}
\usepackage[margin=1in]{geometry}
\usepackage{xcolor}
\usepackage{hyperref}
\usepackage{listings}
\usepackage{amsmath,amssymb}
\usepackage{booktabs}
\usepackage{fancyhdr}
\usepackage{titlesec}
\usepackage{microtype}
\usepackage{tcolorbox}

% --- Color Palette ---
\definecolor{gpxprimary}{RGB}{14, 116, 144}    % Deep Cyan
\definecolor{gpxsecondary}{RGB}{99, 102, 241}  % Indigo
\definecolor{gpxdark}{RGB}{15, 23, 42}         % Slate Dark
\definecolor{gpxbg}{RGB}{248, 250, 252}        % Light Gray Code BG
\definecolor{gpxkeyword}{RGB}{185, 28, 28}     % Crimson
\definecolor{gpxtype}{RGB}{2, 132, 199}        % Cerulean
\definecolor{gpxstring}{RGB}{22, 101, 52}      % Forest Green
\definecolor{gpxcomment}{RGB}{100, 116, 139}   % Slate Muted

% --- Hyperref Setup ---
\hypersetup{
    colorlinks=true,
    linkcolor=gpxprimary,
    filecolor=gpxsecondary,
    urlcolor=gpxprimary,
    citecolor=gpxsecondary,
    pdftitle={GPX Programming Language Specification & Reference Manual},
    pdfauthor={Nishit (@nishit0072e)},
    pdfsubject={Systems Programming Language Specification},
    pdfkeywords={GPX, compiler, systems programming, RISC-V, TAC, debugger, synthesizer}
}

% --- Syntax Highlighting for GPX ---
\lstdefinelanguage{gpx}{
    keywords={let, const, fn, return, if, else, while, for, repeat, when, match, break, continue, struct, union, enum, as, sizeof, null, true, false, in},
    keywordstyle=\color{gpxkeyword}\bfseries,
    keywords=[2]{int, double, float, char, bool, void, string, byte, short, long, *int, *void, *char, List, Node, Stack, Queue, HashMap, BSTNode, Graph},
    keywordstyle=[2]{\color{gpxtype}\bfseries},
    keywords=[3]{print, printf, type, typeof, malloc, free, realloc, calloc, fopen, fclose, fread, fwrite, draw_init, draw_forward, draw_left, draw_right, draw_render_svg, sound_sfx_coin, sound_sfx_laser, sound_play_note, sound_note_to_freq, list_create, list_push, list_pop, list_len},
    keywordstyle=[3]{\color{gpxsecondary}\bfseries},
    sensitive=true,
    comment=[l]{//},
    morecomment=[s]{/*}{*/},
    commentstyle=\color{gpxcomment}\itshape,
    string=[b]",
    morestring=[b]',
    stringstyle=\color{gpxstring}
}

\lstset{
    language=gpx,
    backgroundcolor=\color{gpxbg},
    basicstyle=\ttfamily\small,
    breaklines=true,
    captionpos=b,
    keepspaces=true,
    numbers=left,
    numbersep=8pt,
    numberstyle=\tiny\color{gpxcomment},
    frame=single,
    rulecolor=\color{gray!30},
    showspaces=false,
    showstringspaces=false,
    tabsize=4
}

% --- Headers & Footers ---
\pagestyle{fancy}
\fancyhf{}
\fancyhead[L]{\textcolor{gpxcomment}{\small GPX Language Specification v0.9.1}}
\fancyhead[R]{\textcolor{gpxcomment}{\small \nouppercase{\leftmark}}}
\fancyfoot[C]{\thepage}
\renewcommand{\headrulewidth}{0.4pt}

% --- Title Page ---
\begin{document}

\begin{titlepage}
    \centering
    \vspace*{1.5cm}
    
    {\Huge \bfseries \color{gpxprimary} GPX Programming Language}\\[0.5cm]
    {\LARGE \textbf{Guided Programming eXperience}}\\[0.2cm]
    {\large \textit{Official Language Specification, Architecture \& Reference Manual}}\\[1.5cm]
    
    \rule{0.6\textwidth}{1.5pt}\\[1.5cm]
    
    {\large \textbf{Version 0.9.1}}\\[0.8cm]
    
    \textbf{Lead Architect:}\\
    Nishit (@nishit0072e)\\[0.4cm]
    
    \textbf{Release Date:}\\
    October 2026\\[2cm]
    
    \begin{tcolorbox}[colback=gpxbg,colframe=gpxprimary,arc=3mm,width=0.85\textwidth]
        \centering
        \small
        \textit{"Most people look at compilers and see cold, impenetrable black boxes. We looked at them and asked an audacious question: What if compiling code wasn't a chore, but an exhilarating journey of discovery? What if every developer could peek behind the curtain of modern silicon and feel the raw pulse of the machine?"}
    \end{tcolorbox}
    
    \vfill
    {\footnotesize Copyright \copyright\ 2026 Nishit. All Rights Reserved. Proprietary \& Confidential.}
\end{titlepage}

\tableofcontents
\newpage

% ==============================================================================
\chapter{Introduction \& Design Philosophy}
% ==============================================================================

The \textbf{Guided Programming eXperience (GPX)} is a statically-typed, target-agnostic systems programming language designed for developers, systems engineers, and students who crave bare-metal speed without cognitive obscurity.

\section{Core Tenets}
\begin{itemize}
    \item \textbf{Total Pipeline Introspectability:} Every transformation in GPX---from token streams (\texttt{--tokens}) and Abstract Syntax Trees (\texttt{--ast}) to Three-Address Code (\texttt{--ir}), optimization passes (\texttt{--opt}), and machine code---can be inspected with a single flag.
    \item \textbf{Zero-Overhead Runtime:} GPX executes without garbage collection pauses, heavy runtimes, or hidden virtualization penalties.
    \item \textbf{Dual Audience Design:} GPX seamlessly bridges the gap between educational accessibility (friendly loop constructs, turtle graphics, interactive sound synthesis) and rigorous systems programming (raw pointers, unions, direct memory arenas, and RISC-V assembly generation).
    \item \textbf{Pluggable Multi-Backend Compilation:} Emits high-performance ISO C99 (compiled with GCC) and bare-metal RISC-V RV32I assembly.
\end{itemize}

% ==============================================================================
\chapter{Lexical Structure \& Grammar}
% ==============================================================================

\section{Keywords}
The following tokens are reserved keywords in GPX:
\begin{center}
\begin{tabular}{lllll}
\texttt{let} & \texttt{const} & \texttt{fn} & \texttt{return} & \texttt{if} \\
\texttt{else} & \texttt{while} & \texttt{for} & \texttt{repeat} & \texttt{when} \\
\texttt{break} & \texttt{continue} & \texttt{struct} & \texttt{union} & \texttt{enum} \\
\texttt{as} & \texttt{sizeof} & \texttt{null} & \texttt{true} & \texttt{false} \\
\end{tabular}
\end{center}

\section{Type System}
GPX is statically typed with explicit annotations and local type inference:
\begin{itemize}
    \item \textbf{Primitives:} \texttt{int} (32-bit), \texttt{long} (64-bit), \texttt{short} (16-bit), \texttt{byte} (8-bit), \texttt{char} (8-bit ASCII), \texttt{float} (32-bit IEEE 754), \texttt{double} (64-bit IEEE 754), \texttt{bool}, \texttt{void}, \texttt{string}.
    \item \textbf{Pointers:} Pointers are denoted with an asterisk prefix (e.g., \texttt{*int}, \texttt{*void}, \texttt{*Player}). The address-of operator is \texttt{\&}, and dereferencing is performed via \texttt{*} or arrow access \texttt{->}.
    \item \textbf{Fixed-Size Arrays:} Sized arrays use the syntax \texttt{[T; N]}, with zero-cost indexed access \texttt{arr[i]}.
    \item \textbf{Composite Types:} User-defined records via \texttt{struct}, \texttt{union}, and scoped \texttt{enum}.
\end{itemize}

% ==============================================================================
\chapter{Object-Oriented Programming (OOP) in GPX}
% ==============================================================================

GPX adopts a zero-overhead receiver method syntax modeled after modern systems languages:

\begin{lstlisting}[language=gpx,caption={Receiver Method and Struct Definition}]
struct Vector2D {
    x: double,
    y: double,
}

// Method with pointer receiver
fn (v: *Vector2D) translate(dx: double, dy: double) {
    v->x = v->x + dx;
    v->y = v->y + dy;
}

fn (v: *Vector2D) length_squared() -> double {
    return v->x * v->x + v->y * v->y;
}

fn main() -> int {
    let pos = Vector2D { x: 10.0, y: 20.0 };
    pos.translate(5.0, -2.0);
    print(pos.length_squared());
    return 0;
}
\end{lstlisting}

% ==============================================================================
\chapter{Data Structures \& Algorithms (DSA) Suite}
% ==============================================================================

GPX includes a high-performance, modular standard library for core algorithms and collections:

\begin{description}
    \item[\texttt{<list.gpx>}:] Contiguous growable dynamic array with geometric capacity scaling ($2\times$).
    \item[\texttt{<linkedlist.gpx>}:] Singly and doubly linked node structures for constant-time head/tail insertion.
    \item[\texttt{<stack.gpx>} \& \texttt{<queue.gpx>}:] LIFO stack and FIFO circular ring buffer.
    \item[\texttt{<hashmap.gpx>}:] Key-value associative map with open-addressing and FNV-1a hashing.
    \item[\texttt{<tree.gpx>}:] Self-balancing and binary search trees with recursive traversals.
    \item[\texttt{<graph.gpx>}:] Adjacency-list graphs with Breadth-First (BFS) and Depth-First (DFS) search.
    \item[\texttt{<algorithms.gpx>}:] Quicksort, Mergesort, Insertion sort, and Binary search.
    \item[\texttt{<pool.gpx>}:] O(1) fixed-size block memory pool allocator for deterministic embedded targets.
\end{description}

% ==============================================================================
\chapter{Interactive REPL, Audio Synthesizer \& Debugger}
% ==============================================================================

\section{Interactive REPL (\texttt{gpx --repl})}
The GPX shell provides immediate feedback with type reflection (\texttt{type(x)}), state persistence across inputs, and meta commands (\texttt{.vars}, \texttt{.funcs}, \texttt{.mem}, \texttt{.ir}, \texttt{.ast}, \texttt{.load}).

\section{8-Bit Audio Synthesis (\texttt{<sound.gpx>})}
Direct hardware audio synthesis without external dependencies:
\begin{itemize}
    \item Real-time tones: \texttt{sound\_beep(freq, duration)}
    \item Note notation: \texttt{sound\_play\_note("C4", 250)}
    \item Retro game SFX: \texttt{sound\_sfx\_coin()}, \texttt{sound\_sfx\_laser()}, \texttt{sound\_sfx\_jump()}
    \item Multi-wave synthesis: Square, Sine, Triangle, Sawtooth, and Noise with direct WAV export.
\end{itemize}

\section{Real-Time Memory \& Instruction Inspector (\texttt{gpx --step} / \texttt{gpx-dbg})}
An interactive debugger providing:
\begin{itemize}
    \item Time-travel step rewinding (\texttt{b})
    \item Live Three-Address Code (TAC) instruction cursor
    \item Real-time heap block visualization (allocation id, size, freed status, raw words)
    \item Full Debug Adapter Protocol (DAP) compliance for VS Code and Antigravity IDE.
\end{itemize}

% ==============================================================================
\chapter{Interactive Story \& Adventure Tutorial}
% ==============================================================================

Launched via \texttt{gpx --tutorial}, the adventure mode (*"The Chronicles of the Motherboard"*) features 7 gamified quest chapters:
\begin{enumerate}
    \item \textbf{The Awakening:} Core generator power calculation and integer arithmetic.
    \item \textbf{The Security Gate:} Cryptographic cipher door with boolean branching.
    \item \textbf{The Power Turbine:} Cycling hyperdrive turbine coils via loop accumulation.
    \item \textbf{The Star Beacon:} Projecting 2D turtle vector diagrams into space.
    \item \textbf{Signal of Hope:} Transmitting 8-bit audio homing frequencies.
    \item \textbf{The Autonomous Drone:} Constructing robots with OOP struct receiver methods.
    \item \textbf{The Data Core:} Mainframe stabilization using dynamic arrays (\texttt{<list.gpx>}).
\end{enumerate}

% ==============================================================================
\chapter{CLI Command Reference}
% ==============================================================================

\begin{table}[h]
\centering
\begin{tabular}{llp{7.5cm}}
\toprule
\textbf{Category} & \textbf{Flag / Command} & \textbf{Description} \\
\midrule
Execution & \texttt{--run} & Direct Virtual Machine execution without intermediate files \\
          & \texttt{--repl} & Launches the interactive shell \\
          & \texttt{--tutorial} & Launches the gamified Story / Adventure Mode \\
\midrule
Debugging & \texttt{--step} / \texttt{--debugger} & Launches the interactive Memory Inspector \\
          & \texttt{--dap} & Launches the headless Debug Adapter Protocol server \\
\midrule
Inspection& \texttt{--tokens} & Emits token stream \\
          & \texttt{--ast} & Emits Abstract Syntax Tree hierarchy \\
          & \texttt{--ir} & Emits Three-Address Code (TAC) IR \\
          & \texttt{--opt} & Emits optimized TAC IR \\
          & \texttt{--check} & Semantic analysis and type checking \\
\midrule
Codegen   & \texttt{-o <file.exe>} & Compiles to standalone native binary via GCC \\
          & \texttt{--emit-c} & Transpiles to ISO C99 \\
          & \texttt{--target riscv} & Emits bare-metal RISC-V RV32I assembly \\
\bottomrule
\end{tabular}
\caption{GPX Compiler Command Line Interface}
\end{table}

\end{document}
